\subsection{Nguồn Gốc Lịch Sử}

Giả định về tính tuyến tính gắn liền với sự ra đời của phân tích hồi quy, bắt nguồn từ ba dòng nghiên cứu độc lập ở nửa cuối thế kỷ 19.

\textbf{Ứng dụng đầu tiên từ lý thuyết vật lý: James D.\ Forbes (1857).}
Nhà nghiên cứu Forbes có một lý thuyết cho rằng logarit của áp suất khí quyển $\log(\text{pres})$ có \textit{mối quan hệ tuyến tính} với điểm sôi của nước; quan hệ này có nền tảng vật lý từ phương trình Clausius-Clapeyron. Tài liệu của Forbes được xem là một trong những ví dụ sớm nhất về việc dùng biểu đồ tóm tắt để minh họa hàm trung bình tuyến tính dựa trên mô hình vật lý.

\textbf{Nguồn gốc thuật ngữ ``Hồi quy'': Sir Francis Galton (1877--1886).}
Nhiều ý tưởng nền tảng của hồi quy xuất hiện lần đầu trong công trình của \textbf{Sir Francis Galton (1822--1911)} về sự di truyền đặc tính qua các thế hệ. Vào năm 1877, ông bắt đầu bằng thí nghiệm trên đậu hoa ngọt (sweet peas) để quan sát sự di truyền kích thước hạt. Đến năm 1886, khi biểu diễn hàm trung bình bằng đường thẳng cho dữ liệu chiều cao qua các thế hệ, Galton nhận ra các giá trị cực đoan ở thế hệ trước có xu hướng ``hồi quy'' (\textit{regress}) về phía giá trị trung bình quần thể ở thế hệ sau. Đây là nguồn gốc lịch sử của thuật ngữ \textit{regression}.

\textbf{Sự tiếp nối: Karl Pearson (1893--1903).}
\textbf{Karl Pearson (1857--1936)} mở rộng sang nghiên cứu di truyền học ở người. Trong giai đoạn 1893--1898, ông thu thập dữ liệu chiều cao của các bà mẹ và con gái trưởng thành; bộ dữ liệu này được Pearson và Lee công bố năm 1903 để kiểm tra liệu các bà mẹ cao có xu hướng sinh con gái cao hay không, thông qua phương trình hồi quy tuyến tính.

\textbf{Lưu ý phạm vi lịch sử.} Ba dòng nghiên cứu trên là lịch sử của việc \textit{mô hình hóa quan hệ bằng đường thẳng}. Nền tảng tính toán, tức phương pháp bình phương nhỏ nhất (OLS), được phát triển độc lập bởi Gauss và Legendre ($\sim$1805--1809), và OLS chỉ đảm bảo ước lượng tối ưu (BLUE theo định lý Gauss-Markov) khi dạng hàm được chỉ định đúng.

\textbf{Hệ quả thực tiễn.} Nếu quan hệ thực sự là phi tuyến nhưng ta vẫn dùng mô hình tuyến tính, ước lượng OLS trở nên chệch và không nhất quán; mọi kiểm định và dự báo đều sai.

% -----------------------------------------------------------------------------
